\documentclass{article}
\usepackage[utf8]{inputenc}
\usepackage[T1]{fontenc}
\usepackage{minted}
\usepackage[a4paper, margin=1in]{geometry}
\usepackage{hyperref}
\usepackage{xcolor}

% Configuração para o minted quebrar linhas longas
\setminted{breaklines=true, breakanywhere=true, fontsize=\scriptsize, linenos=true}
\usemintedstyle{vs}

\title{\textbf{Anexo de Códigos: MSR-TCN}}
\author{Pipeline MLOps Completo}
\date{\today}

\begin{document}
\maketitle
\tableofcontents
\newpage
\section{configuracoes.py}
\textbf{Caminho do arquivo:} \texttt{src/configuracoes.py}

\begin{minted}{python}
"""
Módulo de configurações para o projeto MSR-TCN.
Define os segmentos do mercado financeiro, seus respectivos ativos (tickers) e os caminhos globais.
"""
import os
from typing import Dict, List

# Define os 8 segmentos do mercado financeiro utilizados no artigo para avaliação.
# Estes segmentos cobrem ações brasileiras (BlueChips, SmallCaps), Fundos Imobiliários (FIIs), 
# BDRs, Commodities, Mega-Caps dos EUA, Empresas Tradicionais dos EUA e Câmbio.
SEGMENTOS: Dict[str, List[str]] = {
    'BlueChips': [
        'PETR4.SA', 'VALE3.SA', 'ITUB4.SA', 'BBDC4.SA', 'ABEV3.SA', 
        'WEGE3.SA', 'SUZB3.SA', 'RENT3.SA', 'EQTL3.SA', 'RADL3.SA'
    ],
    'SmallCaps': [
        'MGLU3.SA', 'YDUQ3.SA', 'POMO4.SA', 'TEND3.SA', 
        'TASA4.SA', 'FLRY3.SA', 'MRVE3.SA', 'RAPT4.SA'
    ],
    'FIIs': [
        'KNRI11.SA', 'HGLG11.SA', 'MXRF11.SA', 'HGBS11.SA', 'BRCR11.SA',
        'SHPH11.SA', 'HTMX11.SA', 'HGRE11.SA', 'VRTA11.SA', 'KNCR11.SA'
    ],
    'BDRs': [
        'AAPL34.SA', 'MSFT34.SA', 'AMZO34.SA', 'GOGL34.SA', 
        'DISB34.SA', 'TSLA34.SA', 'JNJB34.SA', 'M1TA34.SA'
    ],
    'CommoditiesExp': [
        'GC=F', 'SI=F', 'CL=F', 'BZ=F', 'ZS=F', 
        'ZC=F', 'KC=F', 'LE=F', 'SB=F', 'HG=F',
        'PL=F', 'PA=F', 'CC=F', 'CT=F', 'NG=F'
    ],
    'MegaCapsTech': [
        'AAPL', 'MSFT', 'NVDA', 'GOOGL', 'AMZN', 
        'META', 'TSLA', 'NFLX', 'ADBE', 'CRM', 
        'AMD', 'INTC', 'CSCO', 'QCOM', 'TXN', 
        'AVGO', 'AMAT', 'MU', 'LRCX', 'INTU'
    ],
    'TradicionaisGlobais': [
        'JNJ', 'PG', 'XOM', 'CVX', 'JPM', 
        'V', 'KO', 'PEP', 'WMT', 'MCD', 
        'NKE', 'DIS', 'BA', 'PFE', 'UNH', 
        'HD', 'VZ', 'T', 'MRK', 'ABT'
    ],
    'CambioGlobal': [
        'EURUSD=X', 'JPY=X', 'GBPUSD=X', 'AUDUSD=X', 'USDCAD=X', 
        'USDCHF=X', 'NZDUSD=X', 'EURGBP=X', 'BRL=X'
    ]
}

# Caminhos absolutos para garantir que a execução seja consistente independente do diretório de trabalho.
DIRETORIO_RAIZ = os.path.abspath(os.path.join(os.path.dirname(__file__), '..'))
DIRETORIO_DADOS = os.path.join(DIRETORIO_RAIZ, 'data')
DIRETORIO_CONFIGS = os.path.join(DIRETORIO_RAIZ, 'configs')
DIRETORIO_RESULTADOS = os.path.join(DIRETORIO_RAIZ, 'results')
DIRETORIO_CHECKPOINTS = os.path.join(DIRETORIO_RAIZ, 'checkpoints')

# Garante que os diretórios existam
for diretorio in [DIRETORIO_DADOS, DIRETORIO_CONFIGS, DIRETORIO_RESULTADOS, DIRETORIO_CHECKPOINTS]:
    os.makedirs(diretorio, exist_ok=True)
\end{minted}

\newpage

\section{utilidades.py}
\textbf{Caminho do arquivo:} \texttt{src/utilidades.py}

\begin{minted}{python}
"""
Módulo de utilidades para o projeto MSR-TCN.
Contém funções de perda (loss), carregadores de dados e pipelines de geração de dataset
com Aumentação de Dados (Data Augmentation) integrada para séries temporais.
"""
import os
import torch
import torch.nn as nn
import numpy as np
import pandas as pd
from torch.utils.data import Dataset

class FocalLoss(nn.Module):
    """
    Função Focal Loss para lidar com o desbalanceamento de classes.
    Fornece redimensionamento dinâmico da entropia cruzada com base na confiança da predição.
    """
    def __init__(self, alpha=None, gamma=2):
        super(FocalLoss, self).__init__()
        self.gamma = gamma
        self.alpha = alpha
        
    def forward(self, inputs, targets):
        ce_loss = nn.CrossEntropyLoss(reduction='none')(inputs, targets)
        pt = torch.exp(-ce_loss)
        focal_loss = ((1 - pt) ** self.gamma) * ce_loss
        
        if self.alpha is not None:
            if isinstance(self.alpha, torch.Tensor):
                alpha_t = self.alpha.to(targets.device)[targets]
            else:
                alpha_t = self.alpha[targets]
            focal_loss = alpha_t * focal_loss
            
        return focal_loss.mean()


class SegmentTimeSeriesDataset(Dataset):
    """
    Dataset para carregar dados de séries temporais financeiras segmentadas.
    Suporta Aumentação de Dados em tempo de execução (Jitter, Time-Warp, Window-Slice, MixUp, Decomposition)
    para mitigar o sobreajuste (overfitting) e lidar com o severo desbalanceamento de classes (~90% classe Hold).
    """
    def __init__(self, file_paths: list, window_size: int = 32, is_train: bool = False, 
                 start_date: str = None, end_date: str = None, scaler: dict = None):
        self.X = []
        self.y = []
        self.dates = []
        self.tickers = []
        self.ret_1d = []
        self.is_train = is_train
        self.scaler = scaler if scaler is not None else {}
        
        features_cols = ['Log_Return', 'Volume', 'RSI', 'MACD', 'MACD_Signal', 'ATR']
        
        temp_dfs = []
        for file_path in file_paths:
            if not os.path.exists(file_path): 
                continue
                
            ticker = os.path.basename(file_path).replace('_full.csv', '')
            df = pd.read_csv(file_path, index_col=0)
            df.index = pd.to_datetime(df.index)
            
            if start_date:
                df = df[df.index >= pd.to_datetime(start_date)]
            if end_date:
                df = df[df.index <= pd.to_datetime(end_date)]
                
            if len(df) > window_size:
                # O alvo (target) é o retorno do dia seguinte
                df['Ret_1d'] = df['Close'].pct_change(1).shift(-1) * 100
                temp_dfs.append((df, ticker))
        
        if not temp_dfs:
            self.X, self.y = None, None
            self.da_flags = {}
            return
            
        # Ajusta o padronizador (scaler) usando apenas os dados de treino para evitar vazamento de dados (data leakage)
        if self.is_train:
            combined_df = pd.concat([d[0] for d in temp_dfs])
            for f in features_cols:
                self.scaler[f] = {'mean': combined_df[f].mean(), 'std': combined_df[f].std()}
        
        # Aplica padronização e extrai janelas deslizantes
        for df, ticker in temp_dfs:
            for f in features_cols:
                mean = self.scaler[f]['mean']
                std = self.scaler[f]['std']
                if std != 0:
                    df[f] = (df[f] - mean) / std
                    
            data_X = df[features_cols].values
            data_y = df['Label'].values
            dates_arr = df.index.values
            ret_1d_arr = df['Ret_1d'].fillna(0.0).values
            
            for i in range(len(df) - window_size):
                self.X.append(data_X[i:(i + window_size)])
                self.y.append(data_y[i + window_size])
                
                if not self.is_train:
                    self.dates.append(dates_arr[i + window_size])
                    self.tickers.append(ticker)
                    self.ret_1d.append(ret_1d_arr[i + window_size])
                
        if self.X:
            self.X = np.array(self.X, dtype=np.float32)
            self.y = np.array(self.y, dtype=np.int64)
        else:
            self.X, self.y = None, None
            
        self.da_flags = {}
            
    def __len__(self):
        return len(self.X) if self.X is not None else 0
        
    def set_da_flags(self, flags: dict):
        """Habilita técnicas específicas de Aumentação de Dados."""
        self.da_flags = flags

    def apply_data_augmentation(self, x: np.ndarray, y: int):
        """Aplica as estratégias de aumentação configuradas à janela de amostra."""
        if self.da_flags.get('use_jitter'):
            x += np.random.normal(0, 0.03, x.shape)
            
        if self.da_flags.get('use_warp'):
            warp_factor = np.zeros_like(x)
            for c in range(x.shape[1]):
                warp_factor[:, c] = np.interp(
                    np.linspace(0, 1, len(x)), 
                    [0, 0.5, 1], 
                    np.random.normal(1, 0.1, 3)
                )
            x *= warp_factor
            
        if self.da_flags.get('use_slice'):
            slice_len = int(len(x) * 0.9)
            start_idx = np.random.randint(0, len(x) - slice_len)
            sliced_x = np.zeros_like(x)
            for c in range(x.shape[1]):
                sliced_x[:, c] = np.interp(
                    np.linspace(0, 1, len(x)), 
                    np.linspace(0, 1, slice_len), 
                    x[start_idx : start_idx + slice_len, c]
                )
            x = sliced_x
            
        if self.da_flags.get('use_mixup'):
            idx2 = np.random.randint(0, len(self.X))
            x2 = self.X[idx2].copy()
            lam = np.random.beta(0.5, 0.5)
            x = lam * x + (1 - lam) * x2
            # Label snapping: mantém o rótulo da amostra dominante
            if lam < 0.5: 
                y = self.y[idx2] 
                
        if self.da_flags.get('use_decomp'):
            # Decomposição Temporal: adiciona ruído ao resíduo, preserva a tendência de média móvel
            trend = np.zeros_like(x)
            for c in range(x.shape[1]):
                trend[:, c] = np.convolve(x[:, c], np.ones(3)/3, mode='same')
            x = trend + (x - trend) + np.random.normal(0, 0.05, x.shape)
            
        return x, y

    def __getitem__(self, idx):
        x = self.X[idx].copy()
        y = self.y[idx]
        
        if self.is_train:
            x, y = self.apply_data_augmentation(x, y)
                
        x_tensor = torch.tensor(x, dtype=torch.float32).transpose(0, 1)
        y_tensor = torch.tensor(y, dtype=torch.long)
        return x_tensor, y_tensor
\end{minted}

\newpage

\section{modelos.py}
\textbf{Caminho do arquivo:} \texttt{src/modelos.py}

\begin{minted}{python}
"""
Módulo de modelos para o projeto MSR-TCN.
Contém as implementações da Decomposição Adaptativa em Subbandas (ASD),
Redes Convolucionais Temporais (TCN) e a arquitetura completa do MSR-TCN,
juntamente com modelos base (baselines) para comparação.
"""
import warnings
import torch
import torch.nn as nn
from torch.nn.utils import weight_norm

# Ignorar FutureWarnings relacionadas ao weight_norm em versões mais recentes do PyTorch
warnings.filterwarnings("ignore", category=FutureWarning, message=".*weight_norm.*")

class AsymmetricASD1D(nn.Module):
    """
    Decomposição Adaptativa em Subbandas (ASD) 1D Assimétrica.
    Decompõe a série temporal de entrada em 3 subbandas:
    - A Camada 1 gera Alta Frequência (Ruído) e Baixa Frequência.
    - A Camada 2 decompõe a Baixa Frequência em Média Frequência (Sazonalidade) e Baixa Frequência (Tendência).
    """
    def __init__(self, in_channels: int, filter_size: int = 5):
        super(AsymmetricASD1D, self).__init__()
        
        padding = filter_size // 2
        
        # Camada 1
        self.L1_U = nn.Conv1d(in_channels, in_channels, kernel_size=filter_size, padding=padding, groups=in_channels)
        self.L1_L = nn.Conv1d(in_channels, in_channels, kernel_size=filter_size, padding=padding, groups=in_channels)
        
        # Camada 2 (Opera apenas na saída L1_L)
        self.L2_U = nn.Conv1d(in_channels, in_channels, kernel_size=filter_size, padding=padding, groups=in_channels)
        self.L2_L = nn.Conv1d(in_channels, in_channels, kernel_size=filter_size, padding=padding, groups=in_channels)
        
        # Decimação por 2 após cada filtro
        self.pool = nn.MaxPool1d(kernel_size=2, stride=2)
        
    def forward(self, x: torch.Tensor):
        # x shape: (batch, in_channels, sequence_length)
        
        # Camada 1
        y1_u = self.pool(self.L1_U(x)) # Alta frequência 1 (Ruído)
        y1_l = self.pool(self.L1_L(x)) # Baixa frequência 1
        
        # Camada 2
        y2_u = self.pool(self.L2_U(y1_l)) # Alta frequência 2 (Sazonalidade)
        y2_l = self.pool(self.L2_L(y1_l)) # Baixa frequência 2 (Tendência)
        
        # Sincronizar as dimensões temporais para as 3 subbandas.
        # y1_u sofreu 1 decimação, y2_u e y2_l sofreram 2 decimações.
        # Max-pool em y1_u para sincronizar as dimensões.
        y1_u_sync = self.pool(y1_u)
        
        return y1_u_sync, y2_u, y2_l


class TemporalBlock(nn.Module):
    """
    Bloco Padrão de Rede Convolucional Temporal (TCN) com convoluções causais dilatadas.
    """
    def __init__(self, n_inputs: int, n_outputs: int, kernel_size: int, stride: int, dilation: int, padding: int, dropout: float = 0.2):
        super(TemporalBlock, self).__init__()
        self.pad = nn.ConstantPad1d((padding, 0), 0)
        self.conv1 = weight_norm(nn.Conv1d(n_inputs, n_outputs, kernel_size,
                                           stride=stride, padding=0, dilation=dilation))
        self.relu1 = nn.LeakyReLU(0.1)
        self.dropout1 = nn.Dropout(dropout)

        self.conv2 = weight_norm(nn.Conv1d(n_outputs, n_outputs, kernel_size,
                                           stride=stride, padding=0, dilation=dilation))
        self.relu2 = nn.LeakyReLU(0.1)
        self.dropout2 = nn.Dropout(dropout)

        self.net = nn.Sequential(self.pad, self.conv1, self.relu1, self.dropout1,
                                 self.pad, self.conv2, self.relu2, self.dropout2)
        
        self.downsample = nn.Conv1d(n_inputs, n_outputs, 1) if n_inputs != n_outputs else None
        self.relu = nn.LeakyReLU(0.1)
        self.init_weights()

    def init_weights(self):
        self.conv1.weight.data.normal_(0, 0.01)
        self.conv2.weight.data.normal_(0, 0.01)
        if self.downsample is not None:
            self.downsample.weight.data.normal_(0, 0.01)

    def forward(self, x: torch.Tensor):
        out = self.net(x)
        res = x if self.downsample is None else self.downsample(x)
        return self.relu(out + res)


class SubbandTCN(nn.Module):
    """
    TCN para processar uma única subbanda de forma independente.
    """
    def __init__(self, num_inputs: int, num_channels: list = [16, 32], kernel_size: int = 3, dropout: float = 0.2):
        super(SubbandTCN, self).__init__()
        layers = []
        num_levels = len(num_channels)
        for i in range(num_levels):
            dilation_size = 2 ** i
            in_channels = num_inputs if i == 0 else num_channels[i-1]
            out_channels = num_channels[i]
            layers += [TemporalBlock(in_channels, out_channels, kernel_size, stride=1, dilation=dilation_size,
                                     padding=(kernel_size-1) * dilation_size, dropout=dropout)]

        self.network = nn.Sequential(*layers)
        self.pool = nn.AdaptiveAvgPool1d(1)

    def forward(self, x: torch.Tensor):
        # x shape: (batch, channels, seq_len)
        out = self.network(x)
        out = self.pool(out) # shape: (batch, channels, 1)
        return out.squeeze(-1) # shape: (batch, channels)


class MSRTCN1D(nn.Module):
    """
    MSR-TCN: Rede Convolucional Temporal Residual Multiescala.
    Combina a Decomposição Adaptativa em Subbandas com blocos TCN causais.
    """
    def __init__(self, in_channels: int, seq_len: int, num_classes: int = 3, kernel_size: int = 3, dropout: float = 0.2):
        super(MSRTCN1D, self).__init__()
        self.asd = AsymmetricASD1D(in_channels=in_channels, filter_size=5)
        
        self.tcn_noise = SubbandTCN(in_channels, num_channels=[16, 32], kernel_size=kernel_size, dropout=dropout)
        self.tcn_seasonality = SubbandTCN(in_channels, num_channels=[16, 32], kernel_size=kernel_size, dropout=dropout)
        self.tcn_trend = SubbandTCN(in_channels, num_channels=[16, 32], kernel_size=kernel_size, dropout=dropout)
        
        self.fc = nn.Sequential(
            nn.Linear(96, 128),
            nn.ReLU(),
            nn.Dropout(dropout if dropout else 0.5),
            nn.Linear(128, num_classes)
        )
        
    def forward(self, x: torch.Tensor):
        noise, seasonality, trend = self.asd(x)
        
        feat_noise = self.tcn_noise(noise)
        feat_seasonality = self.tcn_seasonality(seasonality)
        feat_trend = self.tcn_trend(trend)
        
        # Concatenar características de todas as subbandas
        combined = torch.cat([feat_noise, feat_seasonality, feat_trend], dim=1)
        out = self.fc(combined)
        return out


class BaselineTCN1D(nn.Module):
    """
    TCN de Controle (Full-Band) para comparação justa.
    A quantidade de canais é igual à soma dos canais das 3 subbandas
    do MSR-TCN, garantindo capacidade de aprendizado equivalente (contagem de parâmetros).
    """
    def __init__(self, in_channels: int, seq_len: int, num_classes: int = 3, kernel_size: int = 3, dropout: float = 0.2):
        super(BaselineTCN1D, self).__init__()
        
        # MSR-TCN usa 3 subbandas com [16, 32] canais = total de [48, 96]
        self.tcn = SubbandTCN(in_channels, num_channels=[48, 96], kernel_size=kernel_size, dropout=dropout)
        
        self.fc = nn.Sequential(
            nn.Linear(96, 128),
            nn.ReLU(),
            nn.Dropout(dropout if dropout else 0.5),
            nn.Linear(128, num_classes)
        )
        
    def forward(self, x: torch.Tensor):
        feat = self.tcn(x)
        out = self.fc(feat)
        return out


class SubbandCNN(nn.Module):
    """
    CNN Simples para processar cada subbanda independentemente (Modelo de Ablação).
    """
    def __init__(self, in_channels: int, out_channels: int = 16):
        super(SubbandCNN, self).__init__()
        self.conv1 = nn.Conv1d(in_channels, out_channels, kernel_size=3, padding=1)
        self.relu1 = nn.LeakyReLU(0.1)
        self.pool1 = nn.MaxPool1d(2)
        
        self.conv2 = nn.Conv1d(out_channels, out_channels*2, kernel_size=3, padding=1)
        self.relu2 = nn.LeakyReLU(0.1)
        self.pool2 = nn.MaxPool1d(2)
        
        self.flatten = nn.Flatten()
        
    def forward(self, x: torch.Tensor):
        x = self.pool1(self.relu1(self.conv1(x)))
        x = self.pool2(self.relu2(self.conv2(x)))
        return self.flatten(x)


class MSRCNN1D(nn.Module):
    """
    MSR-CNN (Modelo de Ablação).
    Usa ASD 1D para gerar 3 subbandas e processa cada uma com uma CNN isolada.
    """
    def __init__(self, in_channels: int, seq_len: int, num_classes: int = 3):
        super(MSRCNN1D, self).__init__()
        self.asd = AsymmetricASD1D(in_channels=in_channels, filter_size=5)
        
        self.cnn_noise = SubbandCNN(in_channels, out_channels=16)
        self.cnn_seasonality = SubbandCNN(in_channels, out_channels=16)
        self.cnn_trend = SubbandCNN(in_channels, out_channels=16)
        
        final_seq = max(1, seq_len // 16)
        cnn_out_features = 32 * final_seq 
        
        self.fc = nn.Sequential(
            nn.Linear(cnn_out_features * 3, 128), 
            nn.ReLU(),
            nn.Dropout(0.5),
            nn.Linear(128, num_classes)
        )
        
    def forward(self, x: torch.Tensor):
        noise, seasonality, trend = self.asd(x)
        
        feat_noise = self.cnn_noise(noise)
        feat_seasonality = self.cnn_seasonality(seasonality)
        feat_trend = self.cnn_trend(trend)
        
        combined = torch.cat([feat_noise, feat_seasonality, feat_trend], dim=1)
        out = self.fc(combined)
        return out


class BaselineCNN1D(nn.Module):
    """
    CNN Convencional Full-Band (Modelo de Ablação).
    Mesma contagem de parâmetros do MSR-CNN para comparação justa.
    """
    def __init__(self, in_channels: int, seq_len: int, num_classes: int = 3):
        super(BaselineCNN1D, self).__init__()
        
        self.conv1 = nn.Conv1d(in_channels, 48, kernel_size=3, padding=1)
        self.relu1 = nn.LeakyReLU(0.1)
        self.pool1 = nn.MaxPool1d(2)
        
        self.conv2 = nn.Conv1d(48, 96, kernel_size=3, padding=1)
        self.relu2 = nn.LeakyReLU(0.1)
        self.pool2 = nn.MaxPool1d(2)
        
        self.pool_extra = nn.MaxPool1d(4)
        self.flatten = nn.Flatten()
        
        final_seq = max(1, seq_len // 16)
            
        self.fc = nn.Sequential(
            nn.Linear(96 * final_seq, 128),
            nn.ReLU(),
            nn.Dropout(0.5),
            nn.Linear(128, num_classes)
        )
        
    def forward(self, x: torch.Tensor):
        x = self.pool1(self.relu1(self.conv1(x)))
        x = self.pool2(self.relu2(self.conv2(x)))
        x = self.pool_extra(x)
        x = self.flatten(x)
        out = self.fc(x)
        return out
\end{minted}

\newpage

\section{coletar\_dados.py}
\textbf{Caminho do arquivo:} \texttt{src/pipeline\_dados/coletar\_dados.py}

\begin{minted}{python}
"""
Pipeline de extração de dados e engenharia de características (feature engineering) para o MSR-TCN.
Baixa dados históricos via yfinance, calcula indicadores técnicos
e gera os rótulos (labels) com base em uma abordagem de janela deslizante centralizada.
"""
import os
import sys
import numpy as np
import pandas as pd
import yfinance as yf
from ta.momentum import RSIIndicator
from ta.trend import MACD
from ta.volatility import AverageTrueRange

# Garante a importação a partir de src.configuracoes
sys.path.append(os.path.abspath(os.path.join(os.path.dirname(__file__), '../')))
from configuracoes import SEGMENTOS, DIRETORIO_DADOS

def processar_ativo(ticker: str):
    """
    Baixa dados históricos para um determinado ativo (ticker), calcula indicadores técnicos
    e aplica o algoritmo de rotulagem de janela deslizante centralizada (topos e fundos).
    """
    print(f"Baixando dados para {ticker}...")
    df = yf.download(ticker, start='2009-01-01', end='2024-12-20', progress=False)
    
    if df.empty:
        print(f"Falha ao baixar os dados para {ticker}.")
        return
        
    # Achata as colunas MultiIndex, se o yfinance as retornar dessa forma
    if isinstance(df.columns, pd.MultiIndex):
        df.columns = df.columns.droplevel(1)
        
    # 1. Engenharia de Características (Feature Engineering)
    df['Log_Return'] = np.log(df['Close'] / df['Close'].shift(1))
    
    rsi = RSIIndicator(close=df['Close'], window=14)
    df['RSI'] = rsi.rsi()
    
    macd = MACD(close=df['Close'], window_slow=26, window_fast=12, window_sign=9)
    df['MACD'] = macd.macd()
    df['MACD_Signal'] = macd.macd_signal()
    
    atr = AverageTrueRange(high=df['High'], low=df['Low'], close=df['Close'], window=14)
    df['ATR'] = atr.average_true_range()
    
    # 2. Rotulagem de Topo/Fundo (janela centralizada de 11 dias)
    window = 11
    
    # Inicializa o Rótulo (Label) como 0 (Manter/Hold)
    df['Label'] = 0
    
    # Calcula os mínimos e máximos contínuos para a janela centralizada
    # min_periods=window garante que não rotulemos as bordas com janelas incompletas
    rolling_min = df['Close'].rolling(window=window, center=True, min_periods=window).min()
    rolling_max = df['Close'].rolling(window=window, center=True, min_periods=window).max()
    
    # Se o fechamento atual for o mínimo da janela -> 2 (COMPRAR/BUY)
    df.loc[df['Close'] == rolling_min, 'Label'] = 2
    
    # Se o fechamento atual for o máximo da janela -> 1 (VENDER/SELL)
    df.loc[df['Close'] == rolling_max, 'Label'] = 1
    
    df.dropna(inplace=True)
    
    # Os dados são salvos sem padronização (unscaled) para evitar vazamento de dados. 
    # A padronização será aplicada dinamicamente pelo Dataset do PyTorch usando apenas as estatísticas do conjunto de treino.
    nome = ticker.replace('^', '').replace('=', '_')
    caminho_saida = os.path.join(DIRETORIO_DADOS, f"{nome}_full.csv")
    df.to_csv(caminho_saida)
    print(f"Salvo {ticker} em {caminho_saida}")

def main():
    print("Iniciando o pipeline de ingestão de dados...")
    for segmento, ativos in SEGMENTOS.items():
        print(f"\nProcessando Segmento: {segmento}")
        for ativo in ativos:
            processar_ativo(ativo)
    print("\nIngestão de dados concluída com sucesso.")

if __name__ == "__main__":
    main()
\end{minted}

\newpage

\section{01\_otimizar\_hiperparametros.py}
\textbf{Caminho do arquivo:} \texttt{01\_otimizar\_hiperparametros.py}

\begin{minted}{python}
"""
Script de Otimização de Hiperparâmetros para o MSR-TCN.
Executa uma busca em grade (grid search) para encontrar os hiperparâmetros e 
estratégias de Aumentação de Dados (Data Augmentation) ideais para cada um dos 8 segmentos de mercado.
"""
import os
import json
import itertools
import pandas as pd
import torch
import torch.optim as optim
from torch.utils.data import DataLoader
from tqdm import tqdm

import sys
# Garante que os módulos em src possam ser importados
sys.path.append(os.path.join(os.path.dirname(__file__), 'src'))

from configuracoes import SEGMENTOS, DIRETORIO_DADOS, DIRETORIO_CONFIGS
from modelos import MSRTCN1D
from utilidades import SegmentTimeSeriesDataset, FocalLoss

def otimizar_segmento(nome_segmento: str, ativos_segmento: list, device: torch.device):
    print(f"\n[{nome_segmento}] Carregando dados do segmento...")
    arquivos_completos = [os.path.join(DIRETORIO_DADOS, f"{t.replace('^', '').replace('=', '_')}_full.csv") for t in ativos_segmento]
    
    # Utilizamos um período restrito (2013-2014) para uma busca de hiperparâmetros rápida.
    # A avaliação real utiliza a validação Walk-Forward de 2015 a 2024.
    dataset_treino = SegmentTimeSeriesDataset(arquivos_completos, is_train=True, start_date='2013-01-01', end_date='2013-12-31')
    if len(dataset_treino) == 0:
        print(f"[{nome_segmento}] Sem dados de otimização disponíveis. Ignorando.")
        return None
        
    dataset_validacao = SegmentTimeSeriesDataset(arquivos_completos, is_train=False, start_date='2014-01-01', end_date='2014-12-31', scaler=dataset_treino.scaler)
    
    loader_treino = DataLoader(dataset_treino, batch_size=128, shuffle=True)
    loader_validacao = DataLoader(dataset_validacao, batch_size=128, shuffle=False)
        
    # Pré-calcula os pesos alpha para a Focal Loss com base na distribuição de classes
    alvos_treino = torch.tensor(dataset_treino.y, dtype=torch.long)
    contagens = torch.bincount(alvos_treino, minlength=3).float()
    contagens[contagens == 0] = 1.0
    pesos_alpha = len(dataset_treino.y) / (3.0 * contagens)
    pesos_alpha = pesos_alpha.to(device)
        
    grade = {
        'kernel_size': [3, 5],
        'dropout': [0.2],
        'learning_rate': [1e-3, 5e-4],
        'gamma': [2, 3],
        'use_jitter': [True, False],
        'use_warp': [True, False],
        'use_slice': [True, False],
        'use_mixup': [True, False],
        'use_decomp': [True, False]
    }
    
    chaves = grade.keys()
    combinacoes = list(itertools.product(*grade.values()))
    resultados = []
    
    print(f"[{nome_segmento}] Testando {len(combinacoes)} combinações na grade...")
    
    for idx, combo in enumerate(tqdm(combinacoes, desc=f"Grade {nome_segmento}", leave=False)):
        parametros = dict(zip(chaves, combo))
        dataset_treino.set_da_flags(parametros)
        
        modelo = MSRTCN1D(in_channels=6, seq_len=32, kernel_size=parametros['kernel_size'], dropout=parametros['dropout']).to(device)
        criterio = FocalLoss(alpha=pesos_alpha, gamma=parametros['gamma'])
        otimizador = optim.AdamW(modelo.parameters(), lr=parametros['learning_rate'], weight_decay=1e-4)
        
        melhor_perda_val = float('inf')
        
        # 5 épocas para avaliação rápida na busca em grade
        for epoca in range(5):
            modelo.train()
            for X, y in loader_treino:
                X, y = X.to(device), y.to(device)
                otimizador.zero_grad()
                saidas = modelo(X)
                perda = criterio(saidas, y)
                perda.backward()
                otimizador.step()
                
            modelo.eval()
            perda_val_epoca = 0.0
            with torch.no_grad():
                for X, y in loader_validacao:
                    X, y = X.to(device), y.to(device)
                    saidas = modelo(X)
                    perda = criterio(saidas, y)
                    perda_val_epoca += perda.item() * X.size(0)
            
            perda_val_epoca /= len(dataset_validacao)
            if perda_val_epoca < melhor_perda_val: 
                melhor_perda_val = perda_val_epoca
            
        resultados.append({
            **parametros,
            'val_loss': melhor_perda_val 
        })
        
    df_resultados = pd.DataFrame(resultados).sort_values(by='val_loss', ascending=True)
    
    # Salva a melhor configuração
    melhores_parametros = df_resultados.iloc[0].to_dict()
    caminho_json_melhor = os.path.join(DIRETORIO_CONFIGS, f'melhores_parametros_{nome_segmento}.json')
    
    parametros_finais = {
        **melhores_parametros,
        'seq_len': 32,
        'in_channels': 6
    }
    del parametros_finais['val_loss']
    
    with open(caminho_json_melhor, 'w') as f:
        json.dump(parametros_finais, f, indent=4)
        
    print(f"[{nome_segmento}] Melhor configuração salva em: {caminho_json_melhor}")
    return parametros_finais

def main():
    device = torch.device("mps" if torch.backends.mps.is_available() else "cpu")
    print(f"Iniciando Otimização de Hiperparâmetros usando o dispositivo: {device}")
    
    todos_melhores_parametros = []
    
    for nome_segmento, ativos in SEGMENTOS.items():
        res = otimizar_segmento(nome_segmento, ativos, device)
        if res is not None:
            res['Segmento'] = nome_segmento
            todos_melhores_parametros.append(res)
            
    if todos_melhores_parametros:
        df_resumo = pd.DataFrame(todos_melhores_parametros)
        colunas = ['Segmento'] + [c for c in df_resumo.columns if c != 'Segmento']
        df_resumo = df_resumo[colunas]
        
        # Salva o resumo
        diretorio_resultados = os.path.join(os.path.dirname(__file__), 'results')
        os.makedirs(diretorio_resultados, exist_ok=True)
        caminho_csv = os.path.join(diretorio_resultados, 'resumo_otimizacao.csv')
        df_resumo.to_csv(caminho_csv, index=False)
        print(f"\n[Sucesso] Resumo da otimização salvo em: {caminho_csv}")

if __name__ == "__main__":
    main()
\end{minted}

\newpage

\section{02\_treinamento\_walk\_forward.py}
\textbf{Caminho do arquivo:} \texttt{02\_treinamento\_walk\_forward.py}

\begin{minted}{python}
"""
Script de Treinamento de Validação Cruzada Walk-Forward para o MSR-TCN.
Avalia o modelo ao longo de um período de 10 anos (2015-2024) utilizando uma abordagem
de janela em expansão (expanding window) para evitar vazamento de dados.
"""
import os
import json
import pandas as pd
import torch
import torch.optim as optim
from torch.utils.data import DataLoader
from tqdm import tqdm
from sklearn.metrics import accuracy_score, f1_score

import sys
sys.path.append(os.path.abspath(os.path.join(os.path.dirname(__file__), 'src')))

from configuracoes import SEGMENTOS, DIRETORIO_DADOS, DIRETORIO_CONFIGS, DIRETORIO_CHECKPOINTS, DIRETORIO_RESULTADOS
from modelos import MSRTCN1D, BaselineTCN1D
from utilidades import FocalLoss, SegmentTimeSeriesDataset


def treinar_e_avaliar_modelo(modelo, loader_treino, loader_validacao, loader_teste, device, criterio, 
                             epocas=10, caminho_salvar='modelo.pt', lr=1e-3, caminho_plotar=None):
    """
    Loop de treinamento padrão com parada antecipada na validação (salvando o melhor modelo).
    """
    otimizador = optim.AdamW(modelo.parameters(), lr=lr, weight_decay=1e-4)
    melhor_perda_val = float('inf')
    
    perdas_treino = []
    perdas_validacao = []
    
    pbar = tqdm(range(epocas), desc=f"Treinando {os.path.basename(caminho_salvar)}", leave=False)
    for epoca in pbar:
        modelo.train()
        perda_treino_epoca = 0.0
        for X, y in loader_treino:
            X, y = X.to(device), y.to(device)
            otimizador.zero_grad()
            saidas = modelo(X)
            perda = criterio(saidas, y)
            perda.backward()
            otimizador.step()
            perda_treino_epoca += perda.item() * X.size(0)
            
        perda_treino_epoca /= len(loader_treino.dataset)
        perdas_treino.append(perda_treino_epoca)
            
        modelo.eval()
        perda_val_epoca = 0.0
        with torch.no_grad():
            for X, y in loader_validacao:
                X, y = X.to(device), y.to(device)
                saidas = modelo(X)
                perda = criterio(saidas, y)
                perda_val_epoca += perda.item() * X.size(0)
        
        perda_val_epoca /= len(loader_validacao.dataset)
        perdas_validacao.append(perda_val_epoca)
        
        pbar.set_postfix({'T_Loss': f"{perda_treino_epoca:.4f}", 'V_Loss': f"{perda_val_epoca:.4f}"})
        
        if perda_val_epoca < melhor_perda_val:
            melhor_perda_val = perda_val_epoca
            torch.save(modelo.state_dict(), caminho_salvar)
            
    if caminho_plotar:
        import matplotlib.pyplot as plt
        os.makedirs(os.path.dirname(caminho_plotar), exist_ok=True)
        plt.figure(figsize=(10, 6))
        plt.plot(range(1, epocas + 1), perdas_treino, label='Treino', color='#2E86AB', linewidth=2)
        plt.plot(range(1, epocas + 1), perdas_validacao, label='Validação', color='#F24236', linewidth=2, linestyle='--')
        plt.title('Curvas de Aprendizado (Loss vs Épocas)', fontsize=14, fontweight='bold')
        plt.xlabel('Época')
        plt.ylabel('Focal Loss')
        plt.legend()
        plt.grid(True, linestyle='--', alpha=0.6)
        plt.tight_layout()
        plt.savefig(caminho_plotar, dpi=300)
        plt.close()
            
    # Avaliação final no conjunto de Teste
    modelo.load_state_dict(torch.load(caminho_salvar))
    modelo.eval()
    todas_preds, todos_alvos = [], []
    with torch.no_grad():
        for X, y in loader_teste:
            X, y = X.to(device), y.to(device)
            saidas = modelo(X)
            _, previso = torch.max(saidas, 1)
            todas_preds.extend(previso.cpu().numpy())
            todos_alvos.extend(y.cpu().numpy())
            
    acc = accuracy_score(todos_alvos, todas_preds)
    f1 = f1_score(todos_alvos, todas_preds, average='macro')
    return acc, f1, todas_preds, todos_alvos


def treinar_e_avaliar_tcn(nome_modelo, segmento, loader_treino, loader_validacao, loader_teste, device, pesos_alpha=None, melhores_parametros=None):
    print(f"  Treinando {nome_modelo} para {segmento}...")
    
    # Fallback para os parâmetros padrão se as configurações estiverem faltando
    if melhores_parametros is None:
        melhores_parametros = {'kernel_size': 3, 'dropout': 0.2, 'learning_rate': 1e-3}
        
    if nome_modelo == 'MSRTCN':
        modelo = MSRTCN1D(
            in_channels=6, 
            seq_len=32,
            kernel_size=melhores_parametros.get('kernel_size', 3),
            dropout=melhores_parametros.get('dropout', 0.2)
        ).to(device)
    else:
        # TCN Baseline com quantidade equivalente de parâmetros
        modelo = BaselineTCN1D(
            in_channels=6,
            seq_len=32,
            kernel_size=melhores_parametros.get('kernel_size', 3),
            dropout=melhores_parametros.get('dropout', 0.2)
        ).to(device)
    
    criterio = FocalLoss(gamma=2, alpha=pesos_alpha)
    
    caminho_salvar = os.path.join(DIRETORIO_CHECKPOINTS, f"{nome_modelo}_{segmento}.pt")
    caminho_plotar = os.path.join(DIRETORIO_RESULTADOS, 'curvas_aprendizado', f"ca_{nome_modelo}_{segmento}.png")
    
    return treinar_e_avaliar_modelo(
        modelo, loader_treino, loader_validacao, loader_teste, device, criterio, 
        epocas=10, caminho_salvar=caminho_salvar, lr=melhores_parametros.get('learning_rate', 1e-3),
        caminho_plotar=caminho_plotar
    )

def main():
    device = torch.device("mps" if torch.backends.mps.is_available() else "cpu")
    print(f"Iniciando Experimento MSR-TCN Walk-Forward (Dispositivo: {device})")
    
    resultados = []
    
    # No artigo original, 5 dos 8 segmentos mostraram um desempenho significativamente superior
    segmentos_vencedores = ['SmallCaps', 'FIIs', 'MegaCapsTech', 'TradicionaisGlobais', 'CambioGlobal']
    segmentos_alvo = {k: v for k, v in SEGMENTOS.items() if k in segmentos_vencedores}
    
    for segmento, ativos in tqdm(segmentos_alvo.items(), desc="Segmentos"):
        print(f"\n--- Segmento: {segmento} ---")
        arquivos_completos = [os.path.join(DIRETORIO_DADOS, f"{t.replace('^', '').replace('=', '_')}_full.csv") for t in ativos]
        
        caminho_json = os.path.join(DIRETORIO_CONFIGS, f"melhores_parametros_{segmento}.json")
        melhores_parametros = None
        if os.path.exists(caminho_json):
            with open(caminho_json, 'r') as f:
                melhores_parametros = json.load(f)
            print(f"  Parâmetros otimizados carregados: LR={melhores_parametros.get('learning_rate')}, Kernel={melhores_parametros.get('kernel_size')}")
        else:
            print(f"  Aviso: {caminho_json} não encontrado. Usando parâmetros padrão.")
            
        todos_alvos_msrtcn, todas_preds_msrtcn = [], []
        todos_alvos_baseline, todas_preds_baseline = [], []
        
        todas_datas, todos_ativos, todos_ret_1d = [], [], []
        
        # Validação Walk-Forward: 2015 a 2024
        # Treino: [ano_teste - 6] a [ano_teste - 2]
        # Validação: [ano_teste - 1]
        # Teste: [ano_teste]
        for ano_teste in range(2015, 2025):
            print(f"  [Walk-Forward] Ano de Teste: {ano_teste}")
            
            inicio_treino = f"{ano_teste - 6}-01-01"
            fim_treino = f"{ano_teste - 2}-12-31"
            inicio_val = f"{ano_teste - 1}-01-01"
            fim_val = f"{ano_teste - 1}-12-31"
            inicio_teste = f"{ano_teste}-01-01"
            fim_teste = f"{ano_teste}-12-31"
            
            dataset_treino = SegmentTimeSeriesDataset(arquivos_completos, is_train=True, start_date=inicio_treino, end_date=fim_treino)
            if len(dataset_treino) == 0: continue
                
            dataset_val = SegmentTimeSeriesDataset(arquivos_completos, is_train=False, start_date=inicio_val, end_date=fim_val, scaler=dataset_treino.scaler)
            dataset_teste = SegmentTimeSeriesDataset(arquivos_completos, is_train=False, start_date=inicio_teste, end_date=fim_teste, scaler=dataset_treino.scaler)
            
            if len(dataset_val) == 0 or len(dataset_teste) == 0: continue
                
            if melhores_parametros: 
                dataset_treino.set_da_flags(melhores_parametros)
                
            todas_datas.extend(dataset_teste.dates)
            todos_ativos.extend(dataset_teste.tickers)
            todos_ret_1d.extend(dataset_teste.ret_1d)
                
            loader_treino = DataLoader(dataset_treino, batch_size=128, shuffle=True)
            loader_val = DataLoader(dataset_val, batch_size=128, shuffle=False)
            loader_teste = DataLoader(dataset_teste, batch_size=128, shuffle=False)
            
            # Pesos dinâmicos da Focal Loss
            y_tensor = torch.tensor(dataset_treino.y, dtype=torch.long)
            contagens = torch.bincount(y_tensor, minlength=3).float()
            contagens[contagens == 0] = 1.0
            pesos_alpha = len(dataset_treino.y) / (3.0 * contagens)
            pesos_alpha = pesos_alpha.to(device)
            
            acc_m, f1_m, p_m, t_m = treinar_e_avaliar_tcn('MSRTCN', f"{segmento}_{ano_teste}", loader_treino, loader_val, loader_teste, device, pesos_alpha, melhores_parametros)
            acc_b, f1_b, p_b, t_b = treinar_e_avaliar_tcn('BaselineTCN', f"{segmento}_{ano_teste}", loader_treino, loader_val, loader_teste, device, pesos_alpha, melhores_parametros)
            
            todas_preds_msrtcn.extend(p_m)
            todos_alvos_msrtcn.extend(t_m)
            todas_preds_baseline.extend(p_b)
            todos_alvos_baseline.extend(t_b)
            
        if len(todos_alvos_msrtcn) == 0: continue
            
        acc_msrtcn = accuracy_score(todos_alvos_msrtcn, todas_preds_msrtcn)
        f1_msrtcn = f1_score(todos_alvos_msrtcn, todas_preds_msrtcn, average='macro')
        acc_baseline = accuracy_score(todos_alvos_baseline, todas_preds_baseline)
        f1_baseline = f1_score(todos_alvos_baseline, todas_preds_baseline, average='macro')
        
        # Salva as predições para avaliação estatística e análise de portfólio
        df_segmento = pd.DataFrame({
            'Data': todas_datas,
            'Ativo': todos_ativos,
            'Ret_1d': todos_ret_1d,
            'Alvo': todos_alvos_msrtcn,
            'BaselineTCN_Pred': todas_preds_baseline,
            'MSRTCN_Pred': todas_preds_msrtcn
        })
        df_segmento.to_csv(os.path.join(DIRETORIO_RESULTADOS, f'predicoes_{segmento}_wf.csv'), index=False)
        
        resultados.append({
            'Segmento': segmento,
            'BaselineTCN_Acc': acc_baseline, 'BaselineTCN_F1': f1_baseline,
            'MSRTCN_Acc': acc_msrtcn, 'MSRTCN_F1': f1_msrtcn
        })
        
    df_resultados = pd.DataFrame(resultados)
    df_resultados.to_csv(os.path.join(DIRETORIO_RESULTADOS, 'resumo_metricas_walk_forward.csv'), index=False)
    
    print("\nResultados Globais MSR-TCN (Walk-Forward):")
    print(df_resultados)
    print("\nTreinamento Walk-Forward concluído com sucesso!")

if __name__ == '__main__':
    main()
\end{minted}

\newpage

\section{03\_avaliar\_estatisticas.py}
\textbf{Caminho do arquivo:} \texttt{03\_avaliar\_estatisticas.py}

\begin{minted}{python}
"""
Script de Avaliação Estatística para o MSR-TCN.
Realiza o Teste de McNemar e testes de Bootstrap para o F1-Score, avaliando
a significância estatística da superioridade do MSR-TCN sobre a TCN de Controle (Baseline).
"""
import os
import glob
import pandas as pd
import numpy as np
import matplotlib.pyplot as plt
import seaborn as sns
from sklearn.metrics import accuracy_score, f1_score
from statsmodels.stats.contingency_tables import mcnemar

import sys
sys.path.append(os.path.abspath(os.path.join(os.path.dirname(__file__), 'src')))

from configuracoes import DIRETORIO_RESULTADOS

def configurar_estilo_visualizacao():
    """Configura o estilo visual dos gráficos (padrão acadêmico/profissional)."""
    plt.rcParams.update({
        'figure.facecolor': '#FAFAFA',
        'axes.facecolor': '#FAFAFA',
        'axes.edgecolor': '#E0E0E0',
        'axes.labelcolor': '#333333',
        'text.color': '#333333',
        'xtick.color': '#666666',
        'ytick.color': '#666666',
        'grid.color': '#F0F0F0',
        'font.size': 12,
        'font.family': 'sans-serif',
        'axes.spines.top': False,
        'axes.spines.right': False
    })
    sns.set_theme(style='white', context='talk')

def calcular_mcnemar_pareado(alvo, pred_a, pred_b):
    """
    Calcula o p-valor do Teste de McNemar para avaliar se a diferença de precisão
    entre dois modelos é estatisticamente significativa.
    """
    correto_a = (pred_a == alvo).values
    correto_b = (pred_b == alvo).values

    a_errou_b_acertou = sum((~correto_a) & correto_b)
    a_acertou_b_errou = sum(correto_a & (~correto_b))
    ambos_acertaram = sum(correto_a & correto_b)
    ambos_erraram = sum((~correto_a) & (~correto_b))

    tabela = [[ambos_acertaram, a_errou_b_acertou],
              [a_acertou_b_errou, ambos_erraram]]
    
    n_discordantes = a_errou_b_acertou + a_acertou_b_errou
    is_exact = n_discordantes < 25
    resultado = mcnemar(tabela, exact=is_exact, correction=True)
    return resultado.pvalue

def calcular_bootstrap_f1(alvo, pred_a, pred_b, n_iteracoes=1000):
    """
    Calcula a significância (p-valor) via Bootstrap para verificar se o Modelo B (MSR-TCN)
    é sistematicamente superior ao Modelo A (Baseline) na métrica de F1-Score.
    """
    f1_a_orig = f1_score(alvo, pred_a, average='macro')
    f1_b_orig = f1_score(alvo, pred_b, average='macro')
    
    n = len(alvo)
    alvo = np.array(alvo)
    pred_a = np.array(pred_a)
    pred_b = np.array(pred_b)
    
    contador_b_melhor = 0
    for _ in range(n_iteracoes):
        indices = np.random.choice(n, n, replace=True)
        f1_a = f1_score(alvo[indices], pred_a[indices], average='macro', zero_division=0)
        f1_b = f1_score(alvo[indices], pred_b[indices], average='macro', zero_division=0)
        if f1_b > f1_a:
            contador_b_melhor += 1
            
    p_valor = 1.0 - (contador_b_melhor / n_iteracoes)
    return p_valor, f1_a_orig, f1_b_orig

def plotar_p_valores(df_resultados, diretorio_saida):
    """
    Gera um gráfico de barras destacando o grau de certeza (-log10 do p-valor)
    da superioridade do MSR-TCN sobre a TCN de Controle em cada segmento.
    """
    fig, ax = plt.subplots(figsize=(12, 6))
    eps = 1e-300
    df_resultados['-log10(p)'] = -np.log10(df_resultados['p_value_F1'].astype(float) + eps)
    cores = ['#388E3C' if p < 0.05 else '#B0BEC5' for p in df_resultados['p_value_F1']]
    
    barras = sns.barplot(x='Segmento', y='-log10(p)', data=df_resultados, palette=cores, ax=ax)
    
    limiar = -np.log10(0.05)
    ax.axhline(limiar, color='#D32F2F', linestyle='--', label=rf'Significância Mínima ($\alpha=0.05$)', linewidth=2)
    
    ax.set_title('Significância da Evolução no F1-Score: MSR-TCN vs TCN Controle (Bootstrap)', fontsize=18, fontweight='bold', pad=20)
    ax.set_ylabel('Grau de Certeza (-log10 p-value)', fontweight='bold')
    ax.set_xlabel('Segmento Econômico', fontweight='bold')
    plt.xticks(rotation=45)
    ax.legend(frameon=False)
    
    for idx, barra in enumerate(barras.patches):
        pval = df_resultados['p_value_F1'].iloc[idx]
        if pval < 0.05:
            ax.text(barra.get_x() + barra.get_width()/2., barra.get_height() + 0.5,
                    f'p<{pval:.3f}', ha='center', va='bottom', fontsize=9, color='#388E3C', fontweight='bold')

    plt.tight_layout()
    plt.savefig(os.path.join(diretorio_saida, 'significancia_f1_segmentos.png'), dpi=300)
    plt.close()

def main():
    configurar_estilo_visualizacao()
    diretorio_analises = os.path.join(DIRETORIO_RESULTADOS, 'analises_estatisticas')
    os.makedirs(diretorio_analises, exist_ok=True)
    
    # Busca os arquivos de predições gerados pela validação Walk-Forward
    arquivos_predicoes = glob.glob(os.path.join(DIRETORIO_RESULTADOS, 'predicoes_*_wf.csv'))
    
    if not arquivos_predicoes:
        print("Execute o script de Treinamento Walk-Forward (02_) primeiro para gerar as predições.")
        return
        
    resultados = []
    
    for caminho in arquivos_predicoes:
        nome_arquivo = os.path.basename(caminho)
        segmento = nome_arquivo.replace('predicoes_', '').replace('_wf.csv', '')
        
        df = pd.read_csv(caminho)
            
        pval_acc = calcular_mcnemar_pareado(df['Alvo'], df['BaselineTCN_Pred'], df['MSRTCN_Pred'])
        pval_f1, f1_base, f1_msr = calcular_bootstrap_f1(df['Alvo'], df['BaselineTCN_Pred'], df['MSRTCN_Pred'])
        acc_msrtcn = accuracy_score(df['Alvo'], df['MSRTCN_Pred'])
        acc_baseline = accuracy_score(df['Alvo'], df['BaselineTCN_Pred'])
        
        resultados.append({
            'Segmento': segmento,
            'Amostras (N)': len(df),
            'Acurácia TCN Controle': f"{acc_baseline*100:.2f}%",
            'Acurácia MSR-TCN': f"{acc_msrtcn*100:.2f}%",
            'F1 TCN Controle': f"{f1_base:.4f}",
            'F1 MSR-TCN': f"{f1_msr:.4f}",
            'p_value_Acc': pval_acc,
            'p_value_F1': pval_f1,
            'MSR_Venceu_F1_Com_Significancia': 'Sim' if pval_f1 < 0.05 else 'Não'
        })
        
    if resultados:
        df_resultados = pd.DataFrame(resultados)
        caminho_csv = os.path.join(diretorio_analises, 'testes_hipotese_segmentos.csv')
        df_resultados.to_csv(caminho_csv, index=False)
        
        print("\n--- Relatório de Testes de Hipótese ---")
        print(df_resultados.to_string(index=False))
        
        plotar_p_valores(df_resultados, diretorio_analises)
        print(f"\nRelatório de testes estatísticos salvo em {caminho_csv}")

if __name__ == '__main__':
    main()
\end{minted}

\newpage

\section{04\_avaliar\_portfolio.py}
\textbf{Caminho do arquivo:} \texttt{04\_avaliar\_portfolio.py}

\begin{minted}{python}
"""
Script de Avaliação Financeira de Portfólio.
Calcula métricas financeiras institucionais (ROI, Log-Retorno, Sharpe, Sortino, Calmar e MDD)
para simular a aplicação das predições do MSR-TCN na gestão de um fundo de investimento.
"""
import os
import pandas as pd
import numpy as np
import matplotlib.pyplot as plt
import seaborn as sns

import sys
sys.path.append(os.path.abspath(os.path.join(os.path.dirname(__file__), 'src')))
from configuracoes import DIRETORIO_RESULTADOS

DIRETORIO_PORTFOLIO = os.path.join(DIRETORIO_RESULTADOS, 'avaliacao_portfolio')

def configurar_estilo():
    """Configurações visuais institucionais para os gráficos de desempenho."""
    plt.style.use('seaborn-v0_8-whitegrid')
    sns.set_palette("husl")
    plt.rcParams['font.family'] = 'sans-serif'
    plt.rcParams['axes.titlesize'] = 16
    plt.rcParams['axes.labelsize'] = 14
    plt.rcParams['xtick.labelsize'] = 12
    plt.rcParams['ytick.labelsize'] = 12
    plt.rcParams['legend.fontsize'] = 12
    plt.rcParams['figure.titlesize'] = 20

def calcular_sharpe(curva_patrimonio, dias_trade=252):
    """Calcula o Índice de Sharpe anualizado assumindo Taxa Livre de Risco = 0"""
    retornos_diarios = np.diff(curva_patrimonio) / curva_patrimonio[:-1]
    desvio = np.std(retornos_diarios)
    if desvio == 0:
        return 0
    return (np.mean(retornos_diarios) / desvio) * np.sqrt(dias_trade)

def calcular_mdd(curva_patrimonio):
    """Calcula o Maximum Drawdown (MDD) absoluto em %."""
    maximo_acumulado = np.maximum.accumulate(curva_patrimonio)
    maximo_acumulado[maximo_acumulado == 0] = 1 # Proteção contra div/0
    quedas = (curva_patrimonio - maximo_acumulado) / maximo_acumulado
    return np.abs(np.min(quedas)) * 100

def calcular_sortino(curva_patrimonio, dias_trade=252):
    """Calcula o Índice de Sortino."""
    retornos_diarios = np.diff(curva_patrimonio) / curva_patrimonio[:-1]
    retornos_negativos = retornos_diarios[retornos_diarios < 0]
    
    desvio_negativo = np.std(retornos_negativos) if len(retornos_negativos) > 0 else 0
    if desvio_negativo == 0:
        return 0
        
    return (np.mean(retornos_diarios) / desvio_negativo) * np.sqrt(dias_trade)

def calcular_calmar(curva_patrimonio, dias_trade=252):
    """Calcula o Índice de Calmar: CAGR / MDD."""
    anos = len(curva_patrimonio) / dias_trade
    if anos == 0: return 0
    
    cagr = ((curva_patrimonio[-1] / curva_patrimonio[0]) ** (1/anos)) - 1
    mdd = calcular_mdd(curva_patrimonio) / 100.0  
    
    if mdd == 0:
        return 0
    return cagr / mdd

def simular_retornos(predicoes, retornos_1d, capital_inicial=10000.0, custo_transacao=0.0003):
    """
    Simula a curva de capital considerando custos de transação (corretagem/spread).
    Posição: 2=Comprado(Long), 1=Vendido(Flat/Hold em long-only).
    """
    n = len(predicoes)
    retornos_diarios_estrategia = np.zeros(n)
    curva_capital = np.zeros(n + 1)
    curva_capital[0] = capital_inicial
    
    posicao = 0.0  
    
    predicoes_arr = np.array(predicoes)
    ret_arr = np.array(retornos_1d) / 100.0
    
    for i in range(n):
        pred = predicoes_arr[i]
        posicao_anterior = posicao
        
        # Estratégia Long-Only: 2 = Comprar, 1 = Vender (Ficar de fora)
        if pred == 2: posicao = 1.0
        elif pred == 1: posicao = 0.0
        
        custo = custo_transacao if posicao != posicao_anterior else 0.0
        retorno_diario = (posicao * ret_arr[i]) - custo
        retornos_diarios_estrategia[i] = retorno_diario
        curva_capital[i+1] = curva_capital[i] * (1 + retorno_diario)
        
    return retornos_diarios_estrategia, curva_capital[1:]

def plotar_crescimento_portfolio(datas, curva_msr, curva_base, curva_bh, nome_segmento, sufixo_titulo=""):
    fig, ax = plt.subplots(figsize=(14, 7))
    
    ax.plot(datas, curva_bh, label='Mercado (Buy & Hold)', color='#B0BEC5', linestyle='--', linewidth=2)
    ax.plot(datas, curva_base, label='Estratégia TCN Controle', color='#F24236', linewidth=2, alpha=0.8)
    ax.plot(datas, curva_msr, label='Estratégia MSR-TCN', color='#2E86AB', linewidth=3)
    
    ax.yaxis.set_major_formatter(plt.FuncFormatter(lambda x, pos: f'R$ {x:,.0f}'.replace(',', '.')))
    ax.set_title(f'Crescimento de Fundo de Investimento: Portfólio {nome_segmento}\n(Capital Inicial R$ 10.000 por Ativo, Ponderação Igualitária){sufixo_titulo}', fontweight='bold', pad=20)
    ax.set_ylabel('Patrimônio Total do Portfólio', fontweight='bold')
    ax.set_xlabel('Período de Validação', fontweight='bold')
    ax.legend(frameon=True, shadow=True, fancybox=True)
    
    plt.tight_layout()
    plt.savefig(os.path.join(DIRETORIO_PORTFOLIO, f'crescimento_portfolio_{nome_segmento}.png'), dpi=300)
    plt.close()

def plotar_barras_comparacao(df_metricas, sufixo_metrica, string_titulo, string_ylabel, nome_saida, is_percentagem=True):
    plt.figure(figsize=(12, 8))
    
    colunas_plotar = [f'Buy&Hold_{sufixo_metrica}', f'Baseline_{sufixo_metrica}', f'MSRTCN_{sufixo_metrica}']
    
    plot_df = pd.melt(df_metricas, id_vars=['Segmento'], value_vars=colunas_plotar,
                      var_name='Estratégia', value_name='Valor')
    
    plot_df['Estratégia'] = plot_df['Estratégia'].map({
        f'Buy&Hold_{sufixo_metrica}': 'Buy & Hold',
        f'Baseline_{sufixo_metrica}': 'TCN Controle',
        f'MSRTCN_{sufixo_metrica}': 'MSR-TCN'
    })
    
    ax = sns.barplot(x='Segmento', y='Valor', hue='Estratégia', data=plot_df, palette=['#B0BEC5', '#F24236', '#2E86AB'])
    
    plt.title(string_titulo, fontweight='bold', pad=20)
    plt.ylabel(string_ylabel, fontweight='bold')
    plt.xlabel('Segmento', fontweight='bold')
    plt.xticks(rotation=45)
    
    for p in ax.patches:
        height = p.get_height()
        if not np.isnan(height) and height != 0:
            fmt = f'{height:.0f}%' if is_percentagem else f'{height:.2f}'
            ax.annotate(fmt, (p.get_x() + p.get_width() / 2., height),
                        ha='center', va='bottom', fontsize=10, fontweight='bold', xytext=(0, 5),
                        textcoords='offset points')
                        
    plt.legend(title='Estratégia', frameon=True)
    plt.tight_layout()
    plt.savefig(os.path.join(DIRETORIO_PORTFOLIO, nome_saida), dpi=300)
    plt.close()

def main():
    configurar_estilo()
    os.makedirs(DIRETORIO_PORTFOLIO, exist_ok=True)
    print("Gerando Avaliação Financeira Institucional (Log-Retorno, Sharpe, MDD, Sortino, Calmar)...")
    
    segmentos_vencedores = ['SmallCaps', 'FIIs', 'MegaCapsTech', 'TradicionaisGlobais', 'CambioGlobal']
    capital_inicial_por_ativo = 10000.0
    
    todas_metricas = []
    
    for segmento in segmentos_vencedores:
        caminho_csv = os.path.join(DIRETORIO_RESULTADOS, f"predicoes_{segmento}_wf.csv")
        if not os.path.exists(caminho_csv):
            print(f"  [{segmento}] Arquivo {caminho_csv} não encontrado. Execute o script de treino.")
            continue
            
        print(f"  Analisando portfólio {segmento}...")
        df = pd.read_csv(caminho_csv)
        df['Data'] = pd.to_datetime(df['Data'])
        df = df.sort_values(['Ativo', 'Data']).reset_index(drop=True)
        
        ativos = df['Ativo'].unique()
        datas_portfolio = np.sort(df['Data'].unique())
        
        portfolio_msr = np.zeros(len(datas_portfolio))
        portfolio_base = np.zeros(len(datas_portfolio))
        portfolio_bh = np.zeros(len(datas_portfolio))
        
        portfolio_msr_sem_custo = np.zeros(len(datas_portfolio))
        portfolio_base_sem_custo = np.zeros(len(datas_portfolio))
        
        metricas_ativos = []
        
        for ativo in ativos:
            df_ativo = df[df['Ativo'] == ativo].copy()
            df_ativo = df_ativo.set_index('Data').reindex(datas_portfolio).fillna({'Ret_1d': 0, 'MSRTCN_Pred': 0, 'BaselineTCN_Pred': 0})
            
            retorno_bh = df_ativo['Ret_1d'].values / 100.0
            curva_bh = capital_inicial_por_ativo * np.cumprod(1 + retorno_bh)
            
            _, curva_msr = simular_retornos(df_ativo['MSRTCN_Pred'].values, df_ativo['Ret_1d'].values, capital_inicial_por_ativo, 0.0003)
            _, curva_base = simular_retornos(df_ativo['BaselineTCN_Pred'].values, df_ativo['Ret_1d'].values, capital_inicial_por_ativo, 0.0003)
            
            _, curva_msr_sem_custo = simular_retornos(df_ativo['MSRTCN_Pred'].values, df_ativo['Ret_1d'].values, capital_inicial_por_ativo, 0.0)
            _, curva_base_sem_custo = simular_retornos(df_ativo['BaselineTCN_Pred'].values, df_ativo['Ret_1d'].values, capital_inicial_por_ativo, 0.0)
            
            portfolio_msr += curva_msr
            portfolio_base += curva_base
            portfolio_bh += curva_bh
            portfolio_msr_sem_custo += curva_msr_sem_custo
            portfolio_base_sem_custo += curva_base_sem_custo
            
            metricas_ativos.append({
                'Segmento': segmento,
                'Ativo': ativo,
                'Buy&Hold_ROI(%)': ((curva_bh[-1] - capital_inicial_por_ativo) / capital_inicial_por_ativo) * 100,
                'Baseline_ROI(%)': ((curva_base[-1] - capital_inicial_por_ativo) / capital_inicial_por_ativo) * 100,
                'MSRTCN_ROI(%)': ((curva_msr[-1] - capital_inicial_por_ativo) / capital_inicial_por_ativo) * 100,
                'Buy&Hold_Log': np.log(curva_bh[-1] / capital_inicial_por_ativo),
                'Baseline_Log': np.log(max(curva_base[-1], 1e-5) / capital_inicial_por_ativo),
                'MSRTCN_Log': np.log(max(curva_msr[-1], 1e-5) / capital_inicial_por_ativo)
            })
            
        plotar_crescimento_portfolio(datas_portfolio, portfolio_msr, portfolio_base, portfolio_bh, segmento)
        plotar_crescimento_portfolio(datas_portfolio, portfolio_msr_sem_custo, portfolio_base_sem_custo, portfolio_bh, f"{segmento}_sem_custos", "\n[SEM TAXA DE CORRETAGEM]")
        
        df_ativos = pd.DataFrame(metricas_ativos)
        
        todas_metricas.append({
            'Segmento': segmento,
            'Ativos': len(ativos),
            'Buy&Hold_ROI_Mean(%)': df_ativos['Buy&Hold_ROI(%)'].mean(),
            'Baseline_ROI_Mean(%)': df_ativos['Baseline_ROI(%)'].mean(),
            'MSRTCN_ROI_Mean(%)': df_ativos['MSRTCN_ROI(%)'].mean(),
            
            'Buy&Hold_Log_Mean': df_ativos['Buy&Hold_Log'].mean(),
            'Baseline_Log_Mean': df_ativos['Baseline_Log'].mean(),
            'MSRTCN_Log_Mean': df_ativos['MSRTCN_Log'].mean(),
            
            'Buy&Hold_Sharpe': calcular_sharpe(portfolio_bh),
            'Baseline_Sharpe': calcular_sharpe(portfolio_base),
            'MSRTCN_Sharpe': calcular_sharpe(portfolio_msr),
            
            'Buy&Hold_Sortino': calcular_sortino(portfolio_bh),
            'Baseline_Sortino': calcular_sortino(portfolio_base),
            'MSRTCN_Sortino': calcular_sortino(portfolio_msr),
            
            'Buy&Hold_MDD(%)': calcular_mdd(portfolio_bh),
            'Baseline_MDD(%)': calcular_mdd(portfolio_base),
            'MSRTCN_MDD(%)': calcular_mdd(portfolio_msr),
            
            'Buy&Hold_Calmar': calcular_calmar(portfolio_bh),
            'Baseline_Calmar': calcular_calmar(portfolio_base),
            'MSRTCN_Calmar': calcular_calmar(portfolio_msr)
        })
        
        df_ativos.to_csv(os.path.join(DIRETORIO_PORTFOLIO, f'roi_detalhado_{segmento}.csv'), index=False)
        
    if not todas_metricas:
        print("Nenhuma métrica computada. Verifique os arquivos CSV de predição.")
        return

    df_metricas = pd.DataFrame(todas_metricas)
    df_metricas.to_csv(os.path.join(DIRETORIO_PORTFOLIO, 'resumo_segmentos_roi.csv'), index=False)
    
    # Gerar os gráficos de barra
    plotar_barras_comparacao(df_metricas, 'ROI_Mean(%)', 'Retorno Total (Média Aritmética) por Segmento', 'ROI Médio (%)', 'comparacao_segmento_roi_media.png', True)
    plotar_barras_comparacao(df_metricas, 'Log_Mean', 'Log-Retorno Médio (Penalização de Outliers)', 'Média do Log-Retorno', 'comparacao_segmento_roi_log.png', False)
    plotar_barras_comparacao(df_metricas, 'Sharpe', 'Sharpe Ratio Institucional (Retorno/Volatilidade Total)', 'Sharpe Ratio', 'comparacao_segmento_sharpe.png', False)
    plotar_barras_comparacao(df_metricas, 'Sortino', 'Sortino Ratio (Retorno / Volatilidade Negativa)', 'Sortino Ratio', 'comparacao_segmento_sortino.png', False)
    plotar_barras_comparacao(df_metricas, 'Calmar', 'Calmar Ratio (Retorno Anual / MDD)', 'Calmar Ratio', 'comparacao_segmento_calmar.png', False)
    # MDD plot deve ter formatação de porcentagem
    plotar_barras_comparacao(df_metricas, 'MDD(%)', 'Maximum Drawdown (Risco de Ruína) - Menor é Melhor', 'MDD (%)', 'comparacao_segmento_mdd.png', True)
    
    print(f"\n✅ Relatórios Institucionais de Portfólio concluídos: {DIRETORIO_PORTFOLIO}")

if __name__ == '__main__':
    main()
\end{minted}

\newpage

\section{06\_benchmark\_custo\_computacional.py}
\textbf{Caminho do arquivo:} \texttt{06\_benchmark\_custo\_computacional.py}

\begin{minted}{python}
"""
Script de Benchmark de Custo Computacional.
Compara empiricamente e teoricamente o custo computacional (MACs, Parâmetros e Latência)
entre a arquitetura proposta MSR-TCN e o baseline TCN de Controle.
"""
import os
import torch
import time
import thop

import sys
sys.path.append(os.path.abspath(os.path.join(os.path.dirname(__file__), 'src')))

from configuracoes import DIRETORIO_RESULTADOS
from modelos import MSRTCN1D, BaselineTCN1D

def medir_macs_e_parametros(modelo, formato_entrada):
    """Mede Operações de Multiplicação e Acumulação (MACs) e número de parâmetros."""
    modelo.eval()
    modelo.cpu()
    entrada_falsa = torch.randn(formato_entrada)
    macs, parametros = thop.profile(modelo, inputs=(entrada_falsa,), verbose=False)
    return macs, parametros

def medir_tempo_inferencia(modelo_fn, device, formato_entrada, num_execucoes=10000):
    """Mede a latência empírica e a vazão (throughput) de inferência no hardware atual."""
    modelo = modelo_fn().to(device)
    modelo.eval()
    entrada_falsa = torch.randn(formato_entrada).to(device)
    
    # Aquecimento (Warm-up)
    with torch.no_grad():
        for _ in range(100):
            _ = modelo(entrada_falsa)
            
    if device.type == 'cuda':
        torch.cuda.synchronize()
    elif device.type == 'mps':
        torch.mps.synchronize()
        
    tempo_inicio = time.perf_counter()
    with torch.no_grad():
        for _ in range(num_execucoes):
            _ = modelo(entrada_falsa)
            
    if device.type == 'cuda':
        torch.cuda.synchronize()
    elif device.type == 'mps':
        torch.mps.synchronize()
        
    tempo_fim = time.perf_counter()
    
    tempo_medio_ms = ((tempo_fim - tempo_inicio) / num_execucoes) * 1000
    tamanho_lote = formato_entrada[0]
    vazao = tamanho_lote / (tempo_medio_ms / 1000)
    return tempo_medio_ms, vazao

def main():
    device = torch.device("mps" if torch.backends.mps.is_available() else "cpu")
    print(f"Iniciando Benchmark de Eficiência Computacional (Dispositivo: {device})")
    formato_entrada = (128, 6, 32)
    
    # 1. Complexidade Teórica
    macs_base, params_base = medir_macs_e_parametros(BaselineTCN1D(in_channels=6, seq_len=32), formato_entrada)
    macs_msr, params_msr = medir_macs_e_parametros(MSRTCN1D(in_channels=6, seq_len=32), formato_entrada)
    
    # 2. Latência Empírica
    tempo_base, vazao_base = medir_tempo_inferencia(lambda: BaselineTCN1D(in_channels=6, seq_len=32), device, formato_entrada, num_execucoes=5000)
    tempo_msr, vazao_msr = medir_tempo_inferencia(lambda: MSRTCN1D(in_channels=6, seq_len=32), device, formato_entrada, num_execucoes=5000)
    
    relatorio = f"""====================================================
RELATÓRIO DE BENCHMARK COMPUTACIONAL
====================================================

1. Complexidade Teórica (Independente de Hardware)
----------------------------------------------------
[TCN Controle] Parâmetros: {params_base:,.0f} | MACs: {macs_base:,.0f}
[MSR-TCN]      Parâmetros: {params_msr:,.0f} | MACs: {macs_msr:,.0f}

> Redução de Parâmetros: {(1 - params_msr/params_base)*100:.2f}%
> Redução de MACs:       {(1 - macs_msr/macs_base)*100:.2f}%

2. Latência Empírica (Teste de Stress, Lote=128)
----------------------------------------------------
Dispositivo de Execução: {device.type.upper()}

[TCN Controle] Tempo de Inferência: {tempo_base:.3f} ms | Vazão: {vazao_base:,.0f} sequências/seg
[MSR-TCN]      Tempo de Inferência: {tempo_msr:.3f} ms | Vazão: {vazao_msr:,.0f} sequências/seg

> Melhoria de Velocidade: {(tempo_base/tempo_msr - 1)*100:.2f}% mais rápido
====================================================
"""
    print(relatorio)
    
    os.makedirs(DIRETORIO_RESULTADOS, exist_ok=True)
    caminho_salvar = os.path.join(DIRETORIO_RESULTADOS, 'relatorio_benchmark_computacional.txt')
    with open(caminho_salvar, 'w') as f:
        f.write(relatorio)
        
    print(f"✅ Relatório de Benchmark salvo em: {caminho_salvar}")

if __name__ == "__main__":
    main()
\end{minted}

\newpage

\section{05\_gerar\_visualizacoes.py}
\textbf{Caminho do arquivo:} \texttt{05\_gerar\_visualizacoes.py}

\begin{minted}{python}
"""
Script de Geração de Visualizações (Data Storytelling).
Produz gráficos de alta qualidade para explicar os resultados do modelo MSR-TCN,
incluindo Matrizes de Confusão, Boxplots de retornos diários por predição, 
e as curvas de crescimento de capital acumulado.
"""
import os
import glob
import pandas as pd
import numpy as np
import matplotlib.pyplot as plt
import seaborn as sns
from sklearn.metrics import confusion_matrix

import sys
sys.path.append(os.path.abspath(os.path.join(os.path.dirname(__file__), 'src')))
from configuracoes import DIRETORIO_RESULTADOS

DIRETORIO_STORYTELLING = os.path.join(DIRETORIO_RESULTADOS, 'data_storytelling')

def configurar_estilo_storytelling():
    plt.style.use('seaborn-v0_8-white')
    sns.set_palette("husl")
    plt.rcParams['font.family'] = 'sans-serif'
    plt.rcParams['font.sans-serif'] = ['Inter', 'Roboto', 'Arial']

def plotar_matriz_confusao(alvos, predicoes, nome_segmento, caminho_salvar):
    cm = confusion_matrix(alvos, predicoes, labels=[0, 1, 2])
    # Calcula proporção por linha (Rótulo Real)
    soma = cm.sum(axis=1, keepdims=True)
    soma[soma == 0] = 1 # Evita divisão por zero
    cm_pct = cm.astype('float') / soma * 100
    
    fig, ax = plt.subplots(figsize=(8, 6))
    sns.heatmap(cm_pct, annot=True, fmt='.1f', cmap='Blues', 
                xticklabels=['MANTER (0)', 'VENDER (1)', 'COMPRAR (2)'], 
                yticklabels=['MANTER (0)', 'VENDER (1)', 'COMPRAR (2)'],
                cbar_kws={'label': 'Proporção (%)'},
                annot_kws={'size': 14, 'weight': 'bold'}, ax=ax)
                
    ax.set_xlabel('Previsão do MSR-TCN', fontweight='bold')
    ax.set_ylabel('Rótulo Real (Tendência)', fontweight='bold')
    ax.set_title(f'Matriz de Confusão MSR-TCN: {nome_segmento}', fontsize=16, fontweight='bold', pad=20)
    plt.tight_layout()
    plt.savefig(caminho_salvar, dpi=300)
    plt.close()

def plotar_boxplots_predicao(df, nome_segmento, diretorio_saida):
    fig, ax = plt.subplots(figsize=(10, 6))
    
    # Mapear cores: HOLD (Cinza), SELL (Vermelho), BUY (Verde)
    sns.boxplot(x='MSRTCN_Pred', y='Ret_1d', data=df, 
                   palette=['#9E9E9E', '#D32F2F', '#388E3C'], ax=ax)
    
    ax.set_xticklabels(['MANTER (0)', 'VENDER (1)', 'COMPRAR (2)'])
    ax.axhline(0, color='#333333', linestyle='--', linewidth=1.5)
    ax.set_title(f'Retornos Diários Reais por Predição: {nome_segmento}', fontsize=16, fontweight='bold', pad=20)
    ax.set_ylabel('Retorno Real de 1 Dia (%)', fontweight='bold')
    ax.set_xlabel('Decisão Algorítmica do MSR-TCN', fontweight='bold')
    
    plt.tight_layout()
    plt.savefig(os.path.join(diretorio_saida, f'boxplots_retornos_{nome_segmento}.png'), dpi=300)
    plt.close()

def plotar_retornos_acumulados_storytelling(ticker, nome_segmento, df_ticker, diretorio_saida):
    def calcular_retornos_estrategia(predicoes, retornos_1d):
        n = len(predicoes)
        retornos_est = np.zeros(n)
        posicao = 0.0  
        
        predicoes_arr = np.array(predicoes)
        ret_arr = np.array(retornos_1d)
        
        for i in range(n):
            if predicoes_arr[i] == 2: # COMPRAR
                posicao = 1.0
            elif predicoes_arr[i] == 1: # VENDER
                posicao = 0.0
            
            retornos_est[i] = posicao * (ret_arr[i] / 100.0)
            
        return np.cumprod(1 + retornos_est)

    df_ticker = df_ticker.sort_values('Data').reset_index(drop=True)
    datas = pd.to_datetime(df_ticker['Data'])
    
    retornos_bnh = np.cumprod(1 + (df_ticker['Ret_1d'].values / 100.0))
    est_msr = calcular_retornos_estrategia(df_ticker['MSRTCN_Pred'].values, df_ticker['Ret_1d'].values)
    est_base = calcular_retornos_estrategia(df_ticker['BaselineTCN_Pred'].values, df_ticker['Ret_1d'].values)
    
    fig, ax = plt.subplots(figsize=(12, 6))
    
    ax.plot(datas, retornos_bnh, label='Mercado (Buy & Hold)', color='#B0BEC5', linestyle='--', linewidth=2)
    ax.plot(datas, est_base, label='Estratégia TCN Controle', color='#F24236', linewidth=2, alpha=0.8)
    ax.plot(datas, est_msr, label='Estratégia MSR-TCN', color='#2E86AB', linewidth=3)
    
    ax.set_title(f'Crescimento de Capital: {ticker} (Segmento {nome_segmento})', fontsize=18, fontweight='bold', pad=20)
    ax.set_ylabel('Multiplicador do Capital (1.0 = Inicial)', fontweight='bold')
    ax.legend(frameon=False, fontsize=12)
    
    idx_max = np.argmax(est_msr)
    ax.annotate(f'Pico MSR-TCN: {est_msr[idx_max]:.2f}x', 
                xy=(datas.iloc[idx_max], est_msr[idx_max]),
                xytext=(datas.iloc[idx_max], est_msr[idx_max] * 1.1),
                arrowprops=dict(facecolor='black', shrink=0.05, width=1, headwidth=5),
                fontsize=10, fontweight='bold')

    plt.tight_layout()
    plt.savefig(os.path.join(diretorio_saida, f'evolucao_capital_{ticker}.png'), dpi=300)
    plt.close()

def main():
    configurar_estilo_storytelling()
    os.makedirs(DIRETORIO_STORYTELLING, exist_ok=True)
    print("Gerando Data Storytelling (Segmentos Vencedores)...")

    segmentos_vencedores = ['SmallCaps', 'FIIs', 'MegaCapsTech', 'TradicionaisGlobais', 'CambioGlobal']
    
    for segmento in segmentos_vencedores:
        caminho_csv = os.path.join(DIRETORIO_RESULTADOS, f"predicoes_{segmento}_wf.csv")
        if not os.path.exists(caminho_csv):
            print(f"  ⚠ CSV não encontrado para {segmento}. Aguardando treinamento.")
            continue
            
        print(f"Processando Segmento: {segmento}...")
        df = pd.read_csv(caminho_csv)
        
        plotar_matriz_confusao(df['Alvo'].values, df['MSRTCN_Pred'].values, segmento, os.path.join(DIRETORIO_STORYTELLING, f"matriz_confusao_{segmento}.png"))
        plotar_boxplots_predicao(df, segmento, DIRETORIO_STORYTELLING)
        
        # Seleciona 1 ativo representativo para contar a história
        ativos = df['Ativo'].unique()
        if len(ativos) > 0:
            representativo = ativos[0]
            if segmento == 'SmallCaps' and 'YDUQ3.SA' in ativos: representativo = 'YDUQ3.SA'
            if segmento == 'FIIs' and 'MXRF11.SA' in ativos: representativo = 'MXRF11.SA'
            if segmento == 'MegaCapsTech' and 'AAPL' in ativos: representativo = 'AAPL'
            if segmento == 'TradicionaisGlobais' and 'JNJ' in ativos: representativo = 'JNJ'
            if segmento == 'CambioGlobal' and 'EURUSD=X' in ativos: representativo = 'EURUSD=X'
            
            df_ativo = df[df['Ativo'] == representativo].copy()
            plotar_retornos_acumulados_storytelling(representativo, segmento, df_ativo, DIRETORIO_STORYTELLING)

    print(f"\n✅ Relatório de Visualizações (Data Storytelling) concluído com sucesso!")
    print(f"Confira a pasta: {DIRETORIO_STORYTELLING}")

if __name__ == '__main__':
    main()
\end{minted}

\newpage


\end{document}
